% ============================================================
% PACKAGE FILE — package_thesis.tex
% This file centralises all package loading, layout settings,
% custom commands, and bibliography configuration.
% Include it in your main file with: % ============================================================
% PACKAGE FILE — package_thesis.tex
% This file centralises all package loading, layout settings,
% custom commands, and bibliography configuration.
% Include it in your main file with: % ============================================================
% PACKAGE FILE — package_thesis.tex
% This file centralises all package loading, layout settings,
% custom commands, and bibliography configuration.
% Include it in your main file with: % ============================================================
% PACKAGE FILE — package_thesis.tex
% This file centralises all package loading, layout settings,
% custom commands, and bibliography configuration.
% Include it in your main file with: \input{Other_Package/package_thesis}
%
% COMPILATION ORDER:
%   LuaLaTeX -> Biber -> LuaLaTeX -> LuaLaTeX
% ============================================================


% ============================================================
% 1. SPACING AND SECTION TITLE FORMATTING
% ============================================================

\usepackage{setspace}    % Controls line spacing
\usepackage{titlesec}    % Customises section/chapter title layout

\singlespacing           % Default: single spacing
%\onehalfspacing         % Uncomment for 1.5 line spacing

% Vertical space before and after headings:
%   \titlespacing*{<level>}{<left>}{<before>}{<after>}
\titlespacing*{\chapter}     {0pt}{20pt}{8pt}
\titlespacing*{\section}     {0pt}{13pt}{4pt}
\titlespacing*{\subsection}  {0pt}{8pt} {2pt}
\titlespacing*{\subsubsection}{0pt}{4pt}{1pt}
\titlespacing*{\paragraph}   {0pt}{2pt} {1pt}

% Chapter title style: large bold, chapter number followed by title
\titleformat{\chapter}[hang]
  {\normalfont\huge\bfseries}{\thechapter}{13pt}{\Huge}


% ============================================================
% 2. FONT AND LANGUAGE SETTINGS
% ============================================================

\usepackage[LGR,T1]{fontenc}   % T1 encoding (supports Western European + Greek)
\usepackage[utf8]{inputenc}     % UTF-8 input (accented letters, special chars)

% Page geometry — matches B5 paper with binding offset for two-sided printing
\usepackage[b5paper,
            top=2cm, bottom=1.6cm,
            left=2.5cm, right=2cm,
            heightrounded,
            bindingoffset=4mm,
            twoside]{geometry}

% Language support — last language in list is the document default
% Add or remove languages as needed
\usepackage[italian,greek,spanish,english]{babel}


% ============================================================
% 3. PAGE LAYOUT: HEADERS AND FOOTERS
% ============================================================

\usepackage{fancyhdr}        % Custom headers and footers
\fancyhf{}                   % Clear all defaults

% Odd pages (right-hand): section name on the left
\fancyhead[LO]{\nouppercase{\rightmark}}
% Even pages (left-hand): chapter name in italics on the right
\fancyhead[RE]{\textit{\nouppercase{\leftmark}}}
% Centred page number in the footer
\fancyfoot[C]{\thepage}

\renewcommand{\headrulewidth}{0.4pt}  % Thin rule under header
\renewcommand{\footrulewidth}{0pt}    % No rule above footer
\pagestyle{fancy}


% ============================================================
% 4. MATHEMATICS AND CHEMISTRY
% ============================================================

% Core maths packages — safe to keep even if not all are used
\usepackage{amsmath}         % Extended maths environments
\usepackage{amssymb}         % Additional maths symbols
\usepackage{mathtools}       % Extends amsmath (e.g. \coloneqq)
\usepackage{gensymb}         % General symbols: \degree, \celsius, etc.

% Chemistry
\usepackage{chemformula}     % Simple chemical formulas: \ch{H2O}
\usepackage[version=4]{mhchem} % More complex formulas: \ce{Fe2O3}
\usepackage{chemfig}         % Structural chemical diagrams

% Unit formatting
\usepackage{siunitx}         % \SI{10}{\micro\metre}, \ang{45}, etc.

% Maths font (sans-serif style for figures/slides — remove if unwanted)
\usepackage{newtxsf}

% Plotting with pgfplots (e.g. in-document graphs)
\usepackage{pgfplots}
\pgfplotsset{compat=1.17}


% ============================================================
% 5. FIGURES, TABLES, AND FLOATS
% ============================================================

\usepackage{graphicx}        % \includegraphics
\graphicspath{{Figures/}}    % Default search path for images

\usepackage{float}           % [H] placement for figures/tables
\usepackage{wrapfig}         % Text wrapped around figures
\usepackage{subfig}          % Subfigures: (a), (b), etc.
\usepackage{caption}         % Customise caption style
\usepackage{capt-of}         % Captions outside float environments
\usepackage{afterpage}       % Insert full-page floats with \afterpage{}

\usepackage{booktabs}        % Professional table rules (\toprule, \midrule, \bottomrule)
\usepackage{tabularx}        % Tables with auto-width columns (X column type)
\usepackage{longtable}       % Tables spanning multiple pages
\usepackage{array}           % Advanced column formatting in tables
\usepackage{multirow}        % Merge rows in tables
\usepackage{makecell}        % Line breaks inside table cells
\usepackage{colortbl}        % Coloured table cells/rows
\usepackage{lscape}          % Rotate content to landscape (page stays portrait)
\usepackage{pdflscape}       % Rotate entire page to landscape in PDF
\usepackage{rotating}        % Rotate tables or figures


% ============================================================
% 6. TEXT, LAYOUT, AND TYPOGRAPHY
% ============================================================

\usepackage{microtype}       % Improves line breaking and spacing
\usepackage{indentfirst}     % Indent first paragraph of each section
\setlength{\parindent}{1em}  % Paragraph indentation width

\usepackage{ragged2e}        % Better text alignment (\RaggedRight, etc.)
\usepackage{multicol}        % Multi-column layout
\usepackage{changepage}      % Adjust margins locally

\usepackage{epigraph}        % Quotations at chapter openings
\setlength{\epigraphwidth}{0.6\textwidth}

\usepackage{xcolor}          % Custom colours: \definecolor, \textcolor
\usepackage{url}             % Typeset URLs correctly
\usepackage{newunicodechar}  % Map Unicode characters to LaTeX commands
%\newunicodechar{−}{-}       % Example: map Unicode minus to hyphen

\usepackage{textcomp}        % Extra text symbols (e.g. \textsuperscript)
\usepackage{lipsum}          % Dummy text for layout testing: \lipsum[1-3]


% ============================================================
% 7. CODE LISTINGS
% ============================================================

\usepackage{listings}        % Source code with syntax highlighting
\usepackage{matlab-prettifier} % MATLAB-specific highlighting


% ============================================================
% 8. CROSS-REFERENCING AND NAVIGATION
% ============================================================

\usepackage{refstyle}        % Custom cross-reference formats (\figref, \tabref)
\usepackage{bookmark}        % PDF bookmarks for navigation
\usepackage{pdfpages}        % Insert external PDF pages

% Hyperlink configuration (already loaded via \documentclass[hidelinks])
\hypersetup{
    colorlinks = true,
    linkcolor  = black,   % Internal links (TOC, \ref)
    citecolor  = cyan     % Citation links
}


% ============================================================
% 9. LISTS
% ============================================================

\usepackage{enumitem}        % Fine-grained control over list spacing

% Remove extra vertical space in lists
\setlist[itemize] {noitemsep, topsep=0pt}
\setlist[enumerate]{topsep=0pt}


% ============================================================
% 10. BIBLIOGRAPHY
% ============================================================
% Backend: Biber (recommended over BibTeX for full Unicode support)
% Style: numeric-comp — compressed numeric citations [1–3]
% Other common styles: authoryear, alphabetic, apa
%
% To use author-year style instead, replace:
%   style=numeric-comp  ->  style=authoryear
% and recompile with Biber.
% ============================================================

\usepackage{csquotes}  % Required by biblatex for correct quotation marks
\usepackage[backend    = biber,
            style      = numeric-comp,
            giveninits = true,
            url        = false,
            doi        = true,
            eprint     = false,
            natbib     = true,
            sorting    = none]{biblatex}

\addbibresource{references.bib}  % Point to your .bib file

% Clean up bibliography entries: suppress urldate, url, note, language
\DeclareFieldFormat{urldate}{}
\AtEveryBibitem{
    \clearfield{note}
    \clearfield{urldate}
    \clearfield{url}
    \clearlist{language}
}


% ============================================================
% 11. CUSTOM COMMANDS AND ABBREVIATIONS
% ============================================================

% --- Accented uppercase letters (for all-caps contexts) ---
\newcommand*{\egrave}{\MakeUppercase{è}{\space}}
\newcommand*{\eacuta}{\MakeUppercase{é}{\space}}
\newcommand*{\ograve}{\MakeUppercase{ò}{\space}}
\newcommand*{\agrave}{\MakeUppercase{à}{\space}}
\newcommand*{\igrave}{\MakeUppercase{ì}{\space}}
\newcommand*{\ugrave}{\MakeUppercase{ù}{\space}}

% --- Shorthand commands ---
\newcommand*{\cors}{\textit}   % Alias for \textit (italic)
\newcommand*{\sen}{\sin}       % Spanish/Italian: \sen for sine

% --- Custom caption style (small italic) ---
% Usage: place \mycaption before the figure/table environment
\newcommand{\mycaption}{\captionsetup{font={small,it}}}

% --- Chemical abbreviations ---
% Add frequently used compound names here to ensure consistent formatting
\newcommand*{\TiOTwo}{TiO$_2$}
%\newcommand*{\FeOThree}{Fe$_2$O$_3$}   % Example for geochemistry

% --- Styled paragraph title (boxed with horizontal rules) ---
% Usage: \styledtitle{Paragraph heading}
\newcommand{\styledtitle}[1]{%
    \noindent\hrulefill\\[2pt]%
    \par\noindent\hspace*{1cm}\textbf{#1}\par\nobreak%
    \noindent\hrulefill%
}

% --- Landscape page with a centred table ---
% Usage: \landscapetable{<table code>}{<caption text>}
\newcommand{\landscapetable}[2]{%
    \begin{landscape}
        \thispagestyle{empty}
        \begin{center}
            \vspace*{\fill}
            \captionof{table}{#2}
            #1
            \vspace*{\fill}
        \end{center}
    \end{landscape}
}


% ============================================================
% 12. ABSTRACT ENVIRONMENT (decorative)
% ============================================================
% Provides a typographically styled abstract block.
% Usage:
%   \begin{abstract}
%     Text of the abstract...
%   \end{abstract}
% ============================================================

\usepackage{fourier-orns}  % Decorative ornaments (\decofourleft, \decofourright)

\newenvironment{abstract}{%
    \begin{center}
        \vspace{0.6cm}
        \decofourleft\quad{\Large\scshape Abstract}\quad\decofourright
    \end{center}
    \vspace{0.3cm}
    \begin{quotation}
        \small\noindent
}{%
    \end{quotation}
    \vspace{0.5cm}
}

% Suppress the automatic date on the title page
%\date{}

%
% COMPILATION ORDER:
%   LuaLaTeX -> Biber -> LuaLaTeX -> LuaLaTeX
% ============================================================


% ============================================================
% 1. SPACING AND SECTION TITLE FORMATTING
% ============================================================

\usepackage{setspace}    % Controls line spacing
\usepackage{titlesec}    % Customises section/chapter title layout

\singlespacing           % Default: single spacing
%\onehalfspacing         % Uncomment for 1.5 line spacing

% Vertical space before and after headings:
%   \titlespacing*{<level>}{<left>}{<before>}{<after>}
\titlespacing*{\chapter}     {0pt}{20pt}{8pt}
\titlespacing*{\section}     {0pt}{13pt}{4pt}
\titlespacing*{\subsection}  {0pt}{8pt} {2pt}
\titlespacing*{\subsubsection}{0pt}{4pt}{1pt}
\titlespacing*{\paragraph}   {0pt}{2pt} {1pt}

% Chapter title style: large bold, chapter number followed by title
\titleformat{\chapter}[hang]
  {\normalfont\huge\bfseries}{\thechapter}{13pt}{\Huge}


% ============================================================
% 2. FONT AND LANGUAGE SETTINGS
% ============================================================

\usepackage[LGR,T1]{fontenc}   % T1 encoding (supports Western European + Greek)
\usepackage[utf8]{inputenc}     % UTF-8 input (accented letters, special chars)

% Page geometry — matches B5 paper with binding offset for two-sided printing
\usepackage[b5paper,
            top=2cm, bottom=1.6cm,
            left=2.5cm, right=2cm,
            heightrounded,
            bindingoffset=4mm,
            twoside]{geometry}

% Language support — last language in list is the document default
% Add or remove languages as needed
\usepackage[italian,greek,spanish,english]{babel}


% ============================================================
% 3. PAGE LAYOUT: HEADERS AND FOOTERS
% ============================================================

\usepackage{fancyhdr}        % Custom headers and footers
\fancyhf{}                   % Clear all defaults

% Odd pages (right-hand): section name on the left
\fancyhead[LO]{\nouppercase{\rightmark}}
% Even pages (left-hand): chapter name in italics on the right
\fancyhead[RE]{\textit{\nouppercase{\leftmark}}}
% Centred page number in the footer
\fancyfoot[C]{\thepage}

\renewcommand{\headrulewidth}{0.4pt}  % Thin rule under header
\renewcommand{\footrulewidth}{0pt}    % No rule above footer
\pagestyle{fancy}


% ============================================================
% 4. MATHEMATICS AND CHEMISTRY
% ============================================================

% Core maths packages — safe to keep even if not all are used
\usepackage{amsmath}         % Extended maths environments
\usepackage{amssymb}         % Additional maths symbols
\usepackage{mathtools}       % Extends amsmath (e.g. \coloneqq)
\usepackage{gensymb}         % General symbols: \degree, \celsius, etc.

% Chemistry
\usepackage{chemformula}     % Simple chemical formulas: \ch{H2O}
\usepackage[version=4]{mhchem} % More complex formulas: \ce{Fe2O3}
\usepackage{chemfig}         % Structural chemical diagrams

% Unit formatting
\usepackage{siunitx}         % \SI{10}{\micro\metre}, \ang{45}, etc.

% Maths font (sans-serif style for figures/slides — remove if unwanted)
\usepackage{newtxsf}

% Plotting with pgfplots (e.g. in-document graphs)
\usepackage{pgfplots}
\pgfplotsset{compat=1.17}


% ============================================================
% 5. FIGURES, TABLES, AND FLOATS
% ============================================================

\usepackage{graphicx}        % \includegraphics
\graphicspath{{Figures/}}    % Default search path for images

\usepackage{float}           % [H] placement for figures/tables
\usepackage{wrapfig}         % Text wrapped around figures
\usepackage{subfig}          % Subfigures: (a), (b), etc.
\usepackage{caption}         % Customise caption style
\usepackage{capt-of}         % Captions outside float environments
\usepackage{afterpage}       % Insert full-page floats with \afterpage{}

\usepackage{booktabs}        % Professional table rules (\toprule, \midrule, \bottomrule)
\usepackage{tabularx}        % Tables with auto-width columns (X column type)
\usepackage{longtable}       % Tables spanning multiple pages
\usepackage{array}           % Advanced column formatting in tables
\usepackage{multirow}        % Merge rows in tables
\usepackage{makecell}        % Line breaks inside table cells
\usepackage{colortbl}        % Coloured table cells/rows
\usepackage{lscape}          % Rotate content to landscape (page stays portrait)
\usepackage{pdflscape}       % Rotate entire page to landscape in PDF
\usepackage{rotating}        % Rotate tables or figures


% ============================================================
% 6. TEXT, LAYOUT, AND TYPOGRAPHY
% ============================================================

\usepackage{microtype}       % Improves line breaking and spacing
\usepackage{indentfirst}     % Indent first paragraph of each section
\setlength{\parindent}{1em}  % Paragraph indentation width

\usepackage{ragged2e}        % Better text alignment (\RaggedRight, etc.)
\usepackage{multicol}        % Multi-column layout
\usepackage{changepage}      % Adjust margins locally

\usepackage{epigraph}        % Quotations at chapter openings
\setlength{\epigraphwidth}{0.6\textwidth}

\usepackage{xcolor}          % Custom colours: \definecolor, \textcolor
\usepackage{url}             % Typeset URLs correctly
\usepackage{newunicodechar}  % Map Unicode characters to LaTeX commands
%\newunicodechar{−}{-}       % Example: map Unicode minus to hyphen

\usepackage{textcomp}        % Extra text symbols (e.g. \textsuperscript)
\usepackage{lipsum}          % Dummy text for layout testing: \lipsum[1-3]


% ============================================================
% 7. CODE LISTINGS
% ============================================================

\usepackage{listings}        % Source code with syntax highlighting
\usepackage{matlab-prettifier} % MATLAB-specific highlighting


% ============================================================
% 8. CROSS-REFERENCING AND NAVIGATION
% ============================================================

\usepackage{refstyle}        % Custom cross-reference formats (\figref, \tabref)
\usepackage{bookmark}        % PDF bookmarks for navigation
\usepackage{pdfpages}        % Insert external PDF pages

% Hyperlink configuration (already loaded via \documentclass[hidelinks])
\hypersetup{
    colorlinks = true,
    linkcolor  = black,   % Internal links (TOC, \ref)
    citecolor  = cyan     % Citation links
}


% ============================================================
% 9. LISTS
% ============================================================

\usepackage{enumitem}        % Fine-grained control over list spacing

% Remove extra vertical space in lists
\setlist[itemize] {noitemsep, topsep=0pt}
\setlist[enumerate]{topsep=0pt}


% ============================================================
% 10. BIBLIOGRAPHY
% ============================================================
% Backend: Biber (recommended over BibTeX for full Unicode support)
% Style: numeric-comp — compressed numeric citations [1–3]
% Other common styles: authoryear, alphabetic, apa
%
% To use author-year style instead, replace:
%   style=numeric-comp  ->  style=authoryear
% and recompile with Biber.
% ============================================================

\usepackage{csquotes}  % Required by biblatex for correct quotation marks
\usepackage[backend    = biber,
            style      = numeric-comp,
            giveninits = true,
            url        = false,
            doi        = true,
            eprint     = false,
            natbib     = true,
            sorting    = none]{biblatex}

\addbibresource{references.bib}  % Point to your .bib file

% Clean up bibliography entries: suppress urldate, url, note, language
\DeclareFieldFormat{urldate}{}
\AtEveryBibitem{
    \clearfield{note}
    \clearfield{urldate}
    \clearfield{url}
    \clearlist{language}
}


% ============================================================
% 11. CUSTOM COMMANDS AND ABBREVIATIONS
% ============================================================

% --- Accented uppercase letters (for all-caps contexts) ---
\newcommand*{\egrave}{\MakeUppercase{è}{\space}}
\newcommand*{\eacuta}{\MakeUppercase{é}{\space}}
\newcommand*{\ograve}{\MakeUppercase{ò}{\space}}
\newcommand*{\agrave}{\MakeUppercase{à}{\space}}
\newcommand*{\igrave}{\MakeUppercase{ì}{\space}}
\newcommand*{\ugrave}{\MakeUppercase{ù}{\space}}

% --- Shorthand commands ---
\newcommand*{\cors}{\textit}   % Alias for \textit (italic)
\newcommand*{\sen}{\sin}       % Spanish/Italian: \sen for sine

% --- Custom caption style (small italic) ---
% Usage: place \mycaption before the figure/table environment
\newcommand{\mycaption}{\captionsetup{font={small,it}}}

% --- Chemical abbreviations ---
% Add frequently used compound names here to ensure consistent formatting
\newcommand*{\TiOTwo}{TiO$_2$}
%\newcommand*{\FeOThree}{Fe$_2$O$_3$}   % Example for geochemistry

% --- Styled paragraph title (boxed with horizontal rules) ---
% Usage: \styledtitle{Paragraph heading}
\newcommand{\styledtitle}[1]{%
    \noindent\hrulefill\\[2pt]%
    \par\noindent\hspace*{1cm}\textbf{#1}\par\nobreak%
    \noindent\hrulefill%
}

% --- Landscape page with a centred table ---
% Usage: \landscapetable{<table code>}{<caption text>}
\newcommand{\landscapetable}[2]{%
    \begin{landscape}
        \thispagestyle{empty}
        \begin{center}
            \vspace*{\fill}
            \captionof{table}{#2}
            #1
            \vspace*{\fill}
        \end{center}
    \end{landscape}
}


% ============================================================
% 12. ABSTRACT ENVIRONMENT (decorative)
% ============================================================
% Provides a typographically styled abstract block.
% Usage:
%   \begin{abstract}
%     Text of the abstract...
%   \end{abstract}
% ============================================================

\usepackage{fourier-orns}  % Decorative ornaments (\decofourleft, \decofourright)

\newenvironment{abstract}{%
    \begin{center}
        \vspace{0.6cm}
        \decofourleft\quad{\Large\scshape Abstract}\quad\decofourright
    \end{center}
    \vspace{0.3cm}
    \begin{quotation}
        \small\noindent
}{%
    \end{quotation}
    \vspace{0.5cm}
}

% Suppress the automatic date on the title page
%\date{}

%
% COMPILATION ORDER:
%   LuaLaTeX -> Biber -> LuaLaTeX -> LuaLaTeX
% ============================================================


% ============================================================
% 1. SPACING AND SECTION TITLE FORMATTING
% ============================================================

\usepackage{setspace}    % Controls line spacing
\usepackage{titlesec}    % Customises section/chapter title layout

\singlespacing           % Default: single spacing
%\onehalfspacing         % Uncomment for 1.5 line spacing

% Vertical space before and after headings:
%   \titlespacing*{<level>}{<left>}{<before>}{<after>}
\titlespacing*{\chapter}     {0pt}{20pt}{8pt}
\titlespacing*{\section}     {0pt}{13pt}{4pt}
\titlespacing*{\subsection}  {0pt}{8pt} {2pt}
\titlespacing*{\subsubsection}{0pt}{4pt}{1pt}
\titlespacing*{\paragraph}   {0pt}{2pt} {1pt}

% Chapter title style: large bold, chapter number followed by title
\titleformat{\chapter}[hang]
  {\normalfont\huge\bfseries}{\thechapter}{13pt}{\Huge}


% ============================================================
% 2. FONT AND LANGUAGE SETTINGS
% ============================================================

\usepackage[LGR,T1]{fontenc}   % T1 encoding (supports Western European + Greek)
\usepackage[utf8]{inputenc}     % UTF-8 input (accented letters, special chars)

% Page geometry — matches B5 paper with binding offset for two-sided printing
\usepackage[b5paper,
            top=2cm, bottom=1.6cm,
            left=2.5cm, right=2cm,
            heightrounded,
            bindingoffset=4mm,
            twoside]{geometry}

% Language support — last language in list is the document default
% Add or remove languages as needed
\usepackage[italian,greek,spanish,english]{babel}


% ============================================================
% 3. PAGE LAYOUT: HEADERS AND FOOTERS
% ============================================================

\usepackage{fancyhdr}        % Custom headers and footers
\fancyhf{}                   % Clear all defaults

% Odd pages (right-hand): section name on the left
\fancyhead[LO]{\nouppercase{\rightmark}}
% Even pages (left-hand): chapter name in italics on the right
\fancyhead[RE]{\textit{\nouppercase{\leftmark}}}
% Centred page number in the footer
\fancyfoot[C]{\thepage}

\renewcommand{\headrulewidth}{0.4pt}  % Thin rule under header
\renewcommand{\footrulewidth}{0pt}    % No rule above footer
\pagestyle{fancy}


% ============================================================
% 4. MATHEMATICS AND CHEMISTRY
% ============================================================

% Core maths packages — safe to keep even if not all are used
\usepackage{amsmath}         % Extended maths environments
\usepackage{amssymb}         % Additional maths symbols
\usepackage{mathtools}       % Extends amsmath (e.g. \coloneqq)
\usepackage{gensymb}         % General symbols: \degree, \celsius, etc.

% Chemistry
\usepackage{chemformula}     % Simple chemical formulas: \ch{H2O}
\usepackage[version=4]{mhchem} % More complex formulas: \ce{Fe2O3}
\usepackage{chemfig}         % Structural chemical diagrams

% Unit formatting
\usepackage{siunitx}         % \SI{10}{\micro\metre}, \ang{45}, etc.

% Maths font (sans-serif style for figures/slides — remove if unwanted)
\usepackage{newtxsf}

% Plotting with pgfplots (e.g. in-document graphs)
\usepackage{pgfplots}
\pgfplotsset{compat=1.17}


% ============================================================
% 5. FIGURES, TABLES, AND FLOATS
% ============================================================

\usepackage{graphicx}        % \includegraphics
\graphicspath{{Figures/}}    % Default search path for images

\usepackage{float}           % [H] placement for figures/tables
\usepackage{wrapfig}         % Text wrapped around figures
\usepackage{subfig}          % Subfigures: (a), (b), etc.
\usepackage{caption}         % Customise caption style
\usepackage{capt-of}         % Captions outside float environments
\usepackage{afterpage}       % Insert full-page floats with \afterpage{}

\usepackage{booktabs}        % Professional table rules (\toprule, \midrule, \bottomrule)
\usepackage{tabularx}        % Tables with auto-width columns (X column type)
\usepackage{longtable}       % Tables spanning multiple pages
\usepackage{array}           % Advanced column formatting in tables
\usepackage{multirow}        % Merge rows in tables
\usepackage{makecell}        % Line breaks inside table cells
\usepackage{colortbl}        % Coloured table cells/rows
\usepackage{lscape}          % Rotate content to landscape (page stays portrait)
\usepackage{pdflscape}       % Rotate entire page to landscape in PDF
\usepackage{rotating}        % Rotate tables or figures


% ============================================================
% 6. TEXT, LAYOUT, AND TYPOGRAPHY
% ============================================================

\usepackage{microtype}       % Improves line breaking and spacing
\usepackage{indentfirst}     % Indent first paragraph of each section
\setlength{\parindent}{1em}  % Paragraph indentation width

\usepackage{ragged2e}        % Better text alignment (\RaggedRight, etc.)
\usepackage{multicol}        % Multi-column layout
\usepackage{changepage}      % Adjust margins locally

\usepackage{epigraph}        % Quotations at chapter openings
\setlength{\epigraphwidth}{0.6\textwidth}

\usepackage{xcolor}          % Custom colours: \definecolor, \textcolor
\usepackage{url}             % Typeset URLs correctly
\usepackage{newunicodechar}  % Map Unicode characters to LaTeX commands
%\newunicodechar{−}{-}       % Example: map Unicode minus to hyphen

\usepackage{textcomp}        % Extra text symbols (e.g. \textsuperscript)
\usepackage{lipsum}          % Dummy text for layout testing: \lipsum[1-3]


% ============================================================
% 7. CODE LISTINGS
% ============================================================

\usepackage{listings}        % Source code with syntax highlighting
\usepackage{matlab-prettifier} % MATLAB-specific highlighting


% ============================================================
% 8. CROSS-REFERENCING AND NAVIGATION
% ============================================================

\usepackage{refstyle}        % Custom cross-reference formats (\figref, \tabref)
\usepackage{bookmark}        % PDF bookmarks for navigation
\usepackage{pdfpages}        % Insert external PDF pages

% Hyperlink configuration (already loaded via \documentclass[hidelinks])
\hypersetup{
    colorlinks = true,
    linkcolor  = black,   % Internal links (TOC, \ref)
    citecolor  = cyan     % Citation links
}


% ============================================================
% 9. LISTS
% ============================================================

\usepackage{enumitem}        % Fine-grained control over list spacing

% Remove extra vertical space in lists
\setlist[itemize] {noitemsep, topsep=0pt}
\setlist[enumerate]{topsep=0pt}


% ============================================================
% 10. BIBLIOGRAPHY
% ============================================================
% Backend: Biber (recommended over BibTeX for full Unicode support)
% Style: numeric-comp — compressed numeric citations [1–3]
% Other common styles: authoryear, alphabetic, apa
%
% To use author-year style instead, replace:
%   style=numeric-comp  ->  style=authoryear
% and recompile with Biber.
% ============================================================

\usepackage{csquotes}  % Required by biblatex for correct quotation marks
\usepackage[backend    = biber,
            style      = numeric-comp,
            giveninits = true,
            url        = false,
            doi        = true,
            eprint     = false,
            natbib     = true,
            sorting    = none]{biblatex}

\addbibresource{references.bib}  % Point to your .bib file

% Clean up bibliography entries: suppress urldate, url, note, language
\DeclareFieldFormat{urldate}{}
\AtEveryBibitem{
    \clearfield{note}
    \clearfield{urldate}
    \clearfield{url}
    \clearlist{language}
}


% ============================================================
% 11. CUSTOM COMMANDS AND ABBREVIATIONS
% ============================================================

% --- Accented uppercase letters (for all-caps contexts) ---
\newcommand*{\egrave}{\MakeUppercase{è}{\space}}
\newcommand*{\eacuta}{\MakeUppercase{é}{\space}}
\newcommand*{\ograve}{\MakeUppercase{ò}{\space}}
\newcommand*{\agrave}{\MakeUppercase{à}{\space}}
\newcommand*{\igrave}{\MakeUppercase{ì}{\space}}
\newcommand*{\ugrave}{\MakeUppercase{ù}{\space}}

% --- Shorthand commands ---
\newcommand*{\cors}{\textit}   % Alias for \textit (italic)
\newcommand*{\sen}{\sin}       % Spanish/Italian: \sen for sine

% --- Custom caption style (small italic) ---
% Usage: place \mycaption before the figure/table environment
\newcommand{\mycaption}{\captionsetup{font={small,it}}}

% --- Chemical abbreviations ---
% Add frequently used compound names here to ensure consistent formatting
\newcommand*{\TiOTwo}{TiO$_2$}
%\newcommand*{\FeOThree}{Fe$_2$O$_3$}   % Example for geochemistry

% --- Styled paragraph title (boxed with horizontal rules) ---
% Usage: \styledtitle{Paragraph heading}
\newcommand{\styledtitle}[1]{%
    \noindent\hrulefill\\[2pt]%
    \par\noindent\hspace*{1cm}\textbf{#1}\par\nobreak%
    \noindent\hrulefill%
}

% --- Landscape page with a centred table ---
% Usage: \landscapetable{<table code>}{<caption text>}
\newcommand{\landscapetable}[2]{%
    \begin{landscape}
        \thispagestyle{empty}
        \begin{center}
            \vspace*{\fill}
            \captionof{table}{#2}
            #1
            \vspace*{\fill}
        \end{center}
    \end{landscape}
}


% ============================================================
% 12. ABSTRACT ENVIRONMENT (decorative)
% ============================================================
% Provides a typographically styled abstract block.
% Usage:
%   \begin{abstract}
%     Text of the abstract...
%   \end{abstract}
% ============================================================

\usepackage{fourier-orns}  % Decorative ornaments (\decofourleft, \decofourright)

\newenvironment{abstract}{%
    \begin{center}
        \vspace{0.6cm}
        \decofourleft\quad{\Large\scshape Abstract}\quad\decofourright
    \end{center}
    \vspace{0.3cm}
    \begin{quotation}
        \small\noindent
}{%
    \end{quotation}
    \vspace{0.5cm}
}

% Suppress the automatic date on the title page
%\date{}

%
% COMPILATION ORDER:
%   LuaLaTeX -> Biber -> LuaLaTeX -> LuaLaTeX
% ============================================================


% ============================================================
% 1. SPACING AND SECTION TITLE FORMATTING
% ============================================================

\usepackage{setspace}    % Controls line spacing
\usepackage{titlesec}    % Customises section/chapter title layout

\singlespacing           % Default: single spacing
%\onehalfspacing         % Uncomment for 1.5 line spacing

% Vertical space before and after headings:
%   \titlespacing*{<level>}{<left>}{<before>}{<after>}
\titlespacing*{\chapter}     {0pt}{20pt}{8pt}
\titlespacing*{\section}     {0pt}{13pt}{4pt}
\titlespacing*{\subsection}  {0pt}{8pt} {2pt}
\titlespacing*{\subsubsection}{0pt}{4pt}{1pt}
\titlespacing*{\paragraph}   {0pt}{2pt} {1pt}

% Chapter title style: large bold, chapter number followed by title
\titleformat{\chapter}[hang]
  {\normalfont\huge\bfseries}{\thechapter}{13pt}{\Huge}


% ============================================================
% 2. FONT AND LANGUAGE SETTINGS
% ============================================================

\usepackage[LGR,T1]{fontenc}   % T1 encoding (supports Western European + Greek)
\usepackage[utf8]{inputenc}     % UTF-8 input (accented letters, special chars)

% Page geometry — matches B5 paper with binding offset for two-sided printing
\usepackage[b5paper,
            top=2cm, bottom=1.6cm,
            left=2.5cm, right=2cm,
            heightrounded,
            bindingoffset=4mm,
            twoside]{geometry}

% Language support — last language in list is the document default
% Add or remove languages as needed
\usepackage[italian,greek,spanish,english]{babel}


% ============================================================
% 3. PAGE LAYOUT: HEADERS AND FOOTERS
% ============================================================

\usepackage{fancyhdr}        % Custom headers and footers
\fancyhf{}                   % Clear all defaults

% Odd pages (right-hand): section name on the left
\fancyhead[LO]{\nouppercase{\rightmark}}
% Even pages (left-hand): chapter name in italics on the right
\fancyhead[RE]{\textit{\nouppercase{\leftmark}}}
% Centred page number in the footer
\fancyfoot[C]{\thepage}

\renewcommand{\headrulewidth}{0.4pt}  % Thin rule under header
\renewcommand{\footrulewidth}{0pt}    % No rule above footer
\pagestyle{fancy}


% ============================================================
% 4. MATHEMATICS AND CHEMISTRY
% ============================================================

% Core maths packages — safe to keep even if not all are used
\usepackage{amsmath}         % Extended maths environments
\usepackage{amssymb}         % Additional maths symbols
\usepackage{mathtools}       % Extends amsmath (e.g. \coloneqq)
\usepackage{gensymb}         % General symbols: \degree, \celsius, etc.

% Chemistry
\usepackage{chemformula}     % Simple chemical formulas: \ch{H2O}
\usepackage[version=4]{mhchem} % More complex formulas: \ce{Fe2O3}
\usepackage{chemfig}         % Structural chemical diagrams

% Unit formatting
\usepackage{siunitx}         % \SI{10}{\micro\metre}, \ang{45}, etc.

% Maths font (sans-serif style for figures/slides — remove if unwanted)
\usepackage{newtxsf}

% Plotting with pgfplots (e.g. in-document graphs)
\usepackage{pgfplots}
\pgfplotsset{compat=1.17}


% ============================================================
% 5. FIGURES, TABLES, AND FLOATS
% ============================================================

\usepackage{graphicx}        % \includegraphics
\graphicspath{{Figures/}}    % Default search path for images

\usepackage{float}           % [H] placement for figures/tables
\usepackage{wrapfig}         % Text wrapped around figures
\usepackage{subfig}          % Subfigures: (a), (b), etc.
\usepackage{caption}         % Customise caption style
\usepackage{capt-of}         % Captions outside float environments
\usepackage{afterpage}       % Insert full-page floats with \afterpage{}

\usepackage{booktabs}        % Professional table rules (\toprule, \midrule, \bottomrule)
\usepackage{tabularx}        % Tables with auto-width columns (X column type)
\usepackage{longtable}       % Tables spanning multiple pages
\usepackage{array}           % Advanced column formatting in tables
\usepackage{multirow}        % Merge rows in tables
\usepackage{makecell}        % Line breaks inside table cells
\usepackage{colortbl}        % Coloured table cells/rows
\usepackage{lscape}          % Rotate content to landscape (page stays portrait)
\usepackage{pdflscape}       % Rotate entire page to landscape in PDF
\usepackage{rotating}        % Rotate tables or figures


% ============================================================
% 6. TEXT, LAYOUT, AND TYPOGRAPHY
% ============================================================

\usepackage{microtype}       % Improves line breaking and spacing
\usepackage{indentfirst}     % Indent first paragraph of each section
\setlength{\parindent}{1em}  % Paragraph indentation width

\usepackage{ragged2e}        % Better text alignment (\RaggedRight, etc.)
\usepackage{multicol}        % Multi-column layout
\usepackage{changepage}      % Adjust margins locally

\usepackage{epigraph}        % Quotations at chapter openings
\setlength{\epigraphwidth}{0.6\textwidth}

\usepackage{xcolor}          % Custom colours: \definecolor, \textcolor
\usepackage{url}             % Typeset URLs correctly
\usepackage{newunicodechar}  % Map Unicode characters to LaTeX commands
%\newunicodechar{−}{-}       % Example: map Unicode minus to hyphen

\usepackage{textcomp}        % Extra text symbols (e.g. \textsuperscript)
\usepackage{lipsum}          % Dummy text for layout testing: \lipsum[1-3]


% ============================================================
% 7. CODE LISTINGS
% ============================================================

\usepackage{listings}        % Source code with syntax highlighting
\usepackage{matlab-prettifier} % MATLAB-specific highlighting


% ============================================================
% 8. CROSS-REFERENCING AND NAVIGATION
% ============================================================

\usepackage{refstyle}        % Custom cross-reference formats (\figref, \tabref)
\usepackage{bookmark}        % PDF bookmarks for navigation
\usepackage{pdfpages}        % Insert external PDF pages

% Hyperlink configuration (already loaded via \documentclass[hidelinks])
\hypersetup{
    colorlinks = true,
    linkcolor  = black,   % Internal links (TOC, \ref)
    citecolor  = cyan     % Citation links
}


% ============================================================
% 9. LISTS
% ============================================================

\usepackage{enumitem}        % Fine-grained control over list spacing

% Remove extra vertical space in lists
\setlist[itemize] {noitemsep, topsep=0pt}
\setlist[enumerate]{topsep=0pt}


% ============================================================
% 10. BIBLIOGRAPHY
% ============================================================
% Backend: Biber (recommended over BibTeX for full Unicode support)
% Style: numeric-comp — compressed numeric citations [1–3]
% Other common styles: authoryear, alphabetic, apa
%
% To use author-year style instead, replace:
%   style=numeric-comp  ->  style=authoryear
% and recompile with Biber.
% ============================================================

\usepackage{csquotes}  % Required by biblatex for correct quotation marks
\usepackage[backend    = biber,
            style      = numeric-comp,
            giveninits = true,
            url        = false,
            doi        = true,
            eprint     = false,
            natbib     = true,
            sorting    = none]{biblatex}

\addbibresource{references.bib}  % Point to your .bib file

% Clean up bibliography entries: suppress urldate, url, note, language
\DeclareFieldFormat{urldate}{}
\AtEveryBibitem{
    \clearfield{note}
    \clearfield{urldate}
    \clearfield{url}
    \clearlist{language}
}


% ============================================================
% 11. CUSTOM COMMANDS AND ABBREVIATIONS
% ============================================================

% --- Accented uppercase letters (for all-caps contexts) ---
\newcommand*{\egrave}{\MakeUppercase{è}{\space}}
\newcommand*{\eacuta}{\MakeUppercase{é}{\space}}
\newcommand*{\ograve}{\MakeUppercase{ò}{\space}}
\newcommand*{\agrave}{\MakeUppercase{à}{\space}}
\newcommand*{\igrave}{\MakeUppercase{ì}{\space}}
\newcommand*{\ugrave}{\MakeUppercase{ù}{\space}}

% --- Shorthand commands ---
\newcommand*{\cors}{\textit}   % Alias for \textit (italic)
\newcommand*{\sen}{\sin}       % Spanish/Italian: \sen for sine

% --- Custom caption style (small italic) ---
% Usage: place \mycaption before the figure/table environment
\newcommand{\mycaption}{\captionsetup{font={small,it}}}

% --- Chemical abbreviations ---
% Add frequently used compound names here to ensure consistent formatting
\newcommand*{\TiOTwo}{TiO$_2$}
%\newcommand*{\FeOThree}{Fe$_2$O$_3$}   % Example for geochemistry

% --- Styled paragraph title (boxed with horizontal rules) ---
% Usage: \styledtitle{Paragraph heading}
\newcommand{\styledtitle}[1]{%
    \noindent\hrulefill\\[2pt]%
    \par\noindent\hspace*{1cm}\textbf{#1}\par\nobreak%
    \noindent\hrulefill%
}

% --- Landscape page with a centred table ---
% Usage: \landscapetable{<table code>}{<caption text>}
\newcommand{\landscapetable}[2]{%
    \begin{landscape}
        \thispagestyle{empty}
        \begin{center}
            \vspace*{\fill}
            \captionof{table}{#2}
            #1
            \vspace*{\fill}
        \end{center}
    \end{landscape}
}


% ============================================================
% 12. ABSTRACT ENVIRONMENT (decorative)
% ============================================================
% Provides a typographically styled abstract block.
% Usage:
%   \begin{abstract}
%     Text of the abstract...
%   \end{abstract}
% ============================================================

\usepackage{fourier-orns}  % Decorative ornaments (\decofourleft, \decofourright)

\newenvironment{abstract}{%
    \begin{center}
        \vspace{0.6cm}
        \decofourleft\quad{\Large\scshape Abstract}\quad\decofourright
    \end{center}
    \vspace{0.3cm}
    \begin{quotation}
        \small\noindent
}{%
    \end{quotation}
    \vspace{0.5cm}
}

% Suppress the automatic date on the title page
%\date{}
